% ============================================================================
% Draft "Relation to Chen and Maung (2023)" paragraph.
% Drop into the related-literature or discussion section via \input, or copy.
% Requires \cite keys: ChenMaung2023, Pagan1984 (bibtex at end of file).
% Adjust notation (f_t, \varepsilon, \Sigma_e, k, m, T) to match the paper.
% ============================================================================

\subsection{Relation to nonparametric time-varying combination}
\label{subsec:chenmaung}

Our exogenous benchmark is closely related to the time-varying forecast combination
framework of \citet{ChenMaung2023}, who model the combination weights nonparametrically as
smooth functions of rescaled time and establish consistency, an oracle property, and
optimal cross-validated bandwidth selection, including a high-dimensional variant with a
group-SCAD penalty. Their analysis treats the individual forecasts
$f_t=(f_{t1},\dots,f_{td})'$ as regressors and rests on their Assumption~A.3, that the
combination error is a martingale difference sequence,
$\mathbb{E}(\varepsilon_{t+1}\mid I_t)=0$ with $f_t\in I_t$. This predeterminedness and
orthogonality condition plays, in their regression formulation, the same role that the
exogeneity of the forecast errors plays in ours: both condition on the forecasts as given
and require the error to be orthogonal to the information already in the forecasts. Our
common-loading covariance $\Sigma_e=q\iota\iota'+D$ encodes this as a common unforecastable
shock $q\iota\iota'$ plus idiosyncratic noise $D$; the additional restriction that $D$ is
diagonal---mutually uncorrelated idiosyncratic errors---is a structural simplification
beyond A.3 that we use only for the closed-form results.

The assumption is well motivated for genuine external forecasts---for example the Survey of
Professional Forecasters---where it follows from forecast efficiency: an efficient
forecaster has incorporated all available information, so the error around the conditional
mean is orthogonal to the forecast. It is on weaker footing when the forecasts are
themselves estimated from macroeconomic data. In their empirical application the individual
forecasts are constructed from FRED-MD predictors, as are ours; in both cases the forecasts
are \emph{generated regressors} carrying first-stage estimation error and mutually
correlated errors arising from overlapping predictors and a shared estimation sample.
Neither the orthogonality condition A.3 nor the diagonal idiosyncratic structure $D$ is then
delivered by forecast efficiency; both become conditioning assumptions whose validity is an
empirical, specification-dependent matter. Importantly, the asymptotic results in
\citet{ChenMaung2023}---the oracle property, the optimal bandwidth rate, and the limiting
distribution---are derived treating the forecasts as observed, so the first-stage estimation
uncertainty is not propagated. This is the standard generated-regressor caveat
\citep{Pagan1984}, and it is most binding in precisely their high-dimensional case, where
many estimated forecasts enter the selection step at once.

This is the gap our endogenous analysis is designed to address. Rather than condition the
first stage away, we model it. We replace the common-loading idealization with the factor
and banded (Toeplitz) covariance structures that estimated, overlapping forecasts actually
exhibit, and we account for the within-subset estimation cost of fitting $k-1$ weights from
$T$ observations. The latter is what converts the exogenous cube-root tuning rule
$k^\ast\asymp T^{1/3}$ into the square-root rule $k^\ast\asymp T^{1/2}$ in the empirically
relevant regime $m\gtrsim T$: estimating the within-subset weights breaks the first-order
cancellation that drives the cube root, and the Kan--Zhou inflation of the optimal subset
size diverges as $m\to T$. Our exogenous results thus share the assumptions---and the
SPF-strong, FRED-MD-approximate status---of \citet{ChenMaung2023}, while our endogenous
results make explicit the first-stage estimation and the error dependence that their
conditioning, and the forecast-combination literature more broadly, leaves implicit.

% ============================================================================
% BibTeX entries
% ============================================================================
% @article{ChenMaung2023,
%   author  = {Chen, Bin and Maung, Kenwin},
%   title   = {Time-varying forecast combination for high-dimensional data},
%   journal = {Journal of Econometrics},
%   volume  = {237},
%   number  = {2},
%   pages   = {105418},
%   year    = {2023},
%   doi     = {10.1016/j.jeconom.2023.01.024}
% }
%
% @article{Pagan1984,
%   author  = {Pagan, Adrian},
%   title   = {Econometric Issues in the Analysis of Regressions with Generated Regressors},
%   journal = {International Economic Review},
%   volume  = {25},
%   number  = {1},
%   pages   = {221--247},
%   year    = {1984}
% }
